\section{Méthodologie}
\label{sec:methodologie}

Nous suivrons les étapes du descriptif, auxquelles nous ajoutons le CNN et le simulateur. Tous les réglages se feront sur le jeu de validation : le jeu de test reste de côté jusqu'à l'évaluation finale.

\begin{enumerate}
  \item \textbf{Acquisition (notebook 01).} Nous récupérons les annotations et les images utiles, et nous justifions le choix du jeu de données.
  \item \textbf{Exploration (notebook 01).} Nous regardons combien d'images compte chaque classe, si elles sont déséquilibrées et quelle taille font les objets. Nous observons aussi des exemples pour repérer les cas difficiles.
  \item \textbf{Prétraitement (notebook 02).} Nous filtrons, recadrons, redimensionnons et normalisons les images. Le découpage en train, validation et test se fait par image d'origine, pour que des recadrages d'une même photo ne se retrouvent pas dans deux jeux différents.
  \item \textbf{Modèles classiques (notebook 03).} Nous extrayons les descripteurs HOG, puis nous comparons SVM, Random Forest et k-NN à un score plancher.
  \item \textbf{CNN de zéro (notebook 04).} Nous vérifions d'abord que le modèle arrive à mémoriser un petit lot d'images. Ensuite, nous changeons un seul réglage à la fois et nous notons chaque essai.
  \item \textbf{Évaluation (notebook 05).} Nous mesurons accuracy, precision, recall et F1 sur le jeu de test, puis nous cherchons à comprendre les erreurs à l'aide de la matrice de confusion et des images mal classées.
  \item \textbf{Application.} Nous testons le modèle final sur de nouvelles images, puis nous le branchons sur le simulateur.
\end{enumerate}

\textbf{Organisation.} Pour éviter les conflits, une seule personne modifie un notebook à la fois. Avant chaque séance, chacun relit le notebook d'un autre membre, et nous notons nos décisions importantes dans un journal de bord.

\textbf{Principaux risques.} Deux risques nous paraissent les plus sérieux. Si la préparation de COCO prend trop de temps, nous passerons à CIFAR et nous discuterons de sa faible résolution dans nos résultats. Si le temps vient à manquer, nous donnerons la priorité au modèle, et le simulateur restera dans sa version actuelle.

\textbf{Usage de l'IA.} Nous avons préparé le simulateur et la structure des notebooks avec l'aide de Claude Code, et la rédaction de ce dossier a aussi été assistée par cet outil. Le choix des données, le prétraitement, l'entraînement et l'analyse des résultats restent notre travail.
